\section{Discussion}\label{sec:discussion}

Three findings deserve a closer reading. First, taking the memory
contract instead of a window shape as the control variable changes
what a deployment can plan. A contract maps directly onto the question
a serving stack asks, how much context fits a given card, and every
capacity claim in Section~\ref{sec:exp_eff} follows arithmetically
from it, whereas a window shape answers that question only through an
indirect conversion. Organizing scoring, allocation, and eviction
around one contract is what makes the capacity result boring in the
best sense: it is bookkeeping, not extrapolation.

Second, the layer-saturation profile appears to be a property of the
model rather than of the task. The profile of
Figure~\ref{fig:layer_budget} is stable across the LongBench corpus
with a per-layer standard deviation of $0.03$, which suggests that the
allocator can be calibrated once and shared, and that layer-adaptive
budgets are a portable component for other inference stacks, for
instance the paged cache management of vLLM~\cite{kwon2023vllm}.

Third, the errors of eviction are not uniform, but not in the way we
first expected. The absolute gaps to the full cache are largest on
TREC ($1.3$) and TriviaQA ($1.2$), the retrieval-style tasks, and
smallest on the multi-hop MuSiQue ($0.5$); however, the margins over
the strongest baseline are narrowest on the abstractive QMSum ($0.7$
against TOVA) and on MuSiQue ($0.9$, versus $2.4$ on average). The pattern reads as a common failure
shared by every eviction policy: multi-hop evidence receives attention
in bursts separated by long gaps, whatever the scoring rule, so all
methods lose mid-chain entries and converge to similar scores, while
on retrieval-style tasks the policies differ most and SAGE's
refresh-on-revisit dynamics pay off. The score trace below makes the
mechanism visible: intermediate hops hold their scores through the
gaps between questions, the decay erodes them under every policy we
logged, and eviction removes exactly the entries the later hops need
to revisit. Query-aware score compensation, or a task-level protection
set seeded from the question, is the natural next step, and the trace
analysis indicates it should target the mid-chain evidence rather
than the extremes of the document.

\section{Limitations}\label{sec:limitations}

Eviction is lossy by construction: once an entry leaves the cache it
cannot be recovered, and below a $10\%$ budget the gap to the full
cache is no longer negligible, $0.80$ perplexity on WikiText-2 for the
7B model. Stacking SAGE with cache
quantization~\cite{liu2024kivi} is
the obvious remedy, since the two act on orthogonal dimensions, and we
have not yet evaluated the combination. The evaluation covers three
dense decoder models of 7B and 13B scale and English benchmarks only;
the task-level, retrieval, and generation evaluations beyond perplexity
run on the 13B model alone, and recent baselines with fixed pyramidal
layer budgets or query-aware selection~\cite{zhang2024pyramidkv} are
cited but not compared head-to-head, so the measured margins describe
the 2023 generation of eviction methods. Mixture-of-experts
architectures, grouped-query attention, multilingual long-context
behavior, and encoder-decoder models remain untested, so the claim of
applicability to any decoder-only model is a design statement rather
than a measured one. The throughput
gains depend on the attention kernel exposing a gather path; our
implementation sits on top of FlashAttention-2 and would benefit from
a native sparse-read kernel, which raises the ceiling rather than the
measured floor. The scoring implementation approximates the tail of
the attention distribution with a shared constant, which introduces a
bounded bias, below $0.02$ perplexity at $m = 64$ in our measurements,
whose formal analysis we leave open. Finally, the bucketed heap trades
strict score ordering for amortized constant eviction time, so the
policy is near-greedy rather than exactly greedy on the score; the
ablation of Table~\ref{tab:gate} already prices this approximation as
part of the framework.

\section{Conclusion}\label{sec:conclusion}

SAGE organizes training-free KV cache eviction as three coupled
mechanisms under a single global memory contract: decayed
attention-confidence scoring, layer-adaptive budget allocation, and a
dual-reservoir cache with a protected recent window. At a $20\%$
contract, three instruction-tuned models retain $98.0\%$ of the
LongBench average and up to $94.8\%$ needle-retrieval accuracy at 128K
tokens, while the cache shrinks by a factor of $4.4$ and decode
throughput rises to $2.02\times$ at the longest context we test. The
ablation attributes an identifiable share of the quality gap to each
component, the sensitivity analysis places the defaults on flat
plateaus, and the layer-profile analysis shows that the budgets follow
a stable model property rather than a task artifact. We release code,
configurations, and the complete experimental record with the paper.

\paragraph{Deployment guidance.}
The results translate into a simple decision rule for practitioners.
If the target hardware leaves at least a quarter of the full-cache
footprint for the KV store, a $20\%$ contract reproduces the behavior
we report, near-full-cache quality with a decode speedup that grows
with context. Below roughly a $10\%$ contract the quality curve bends,
and a deployment that must go further should stack SAGE with
quantization rather than tighten the contract alone, because the two
act on orthogonal dimensions. Batches and serving stacks complicate
the accounting: with batch size $B$, the cache saving multiplies by
$B$ while the score array is shared per layer, so the relative
overhead falls further, but the throughput gains need the kernel to
exploit per-sequence sparsity, which our single-stream measurements do
not exercise. When the latency of the first token dominates, the
$5.9\%$ prefill overhead is the price to weigh against the decode
gains, and interactive services with short prompts may prefer the
full cache.
